\subsection{Steering recipe}
For layer $L$ and MLP neuron $n$, we add a constant $\delta$ to the neuron's activation at the
$\mathrm{down\_proj}$ input at every token position:
\begin{equation}
  \delta \;=\; \mathrm{dose}\times
  \frac{\operatorname{median}_{t\ge 1}\,\lVert h^{(L)}_t\rVert}{\lVert \Wdown[:,n]\rVert},
\end{equation}
where $h^{(L)}_t$ is the residual stream leaving layer $L$ on the unsteered model, measured once per layer
over the six prompts of the original demonstration (Appendix~\ref{app:calib} lists all 64 values; layer 34
gives $93.30$). Because $\mathrm{down\_proj}$ is linear this equals adding $\delta\,\Wdown[:,n]$ to the MLP
output, which is how we implement it (\texttt{AddMLP} in \texttt{hunt/common.py}). A dose of $1$ is a push as
long as the typical residual vector. The implementation reproduces the reference \texttt{steer.py} text for
\neuron{34}{3838} at $0.75$ token for token. All generation is greedy with thinking disabled, on
\texttt{mlx-community/Qwen3.8-27B-8bit}.

\subsection{Candidate ranking}\label{sec:ranking}
For each of 13 target personas (joyful, sad, fearful, evil, in love, pirate, poetic, paranoid, childlike,
mystical, bored, sycophantic, grandiose) plus a neutral condition, the unsteered model answers the 32 TRAIN
prompts under a one-sentence system prompt describing the persona (Appendix~\ref{app:judge}; 120 tokens).
Each answer is then fed back teacher-forced after the plain user prompt, with no system prompt, and we record
every MLP neuron's mean activation over the answer tokens and the mean residual stream leaving each layer. The
persona direction is $d_{p}^{(L)} = \bar h_{p}^{(L)} - \bar h_{\neg p}^{(L)}$, the persona's mean residual
minus the mean over the other thirteen conditions, averaged over prompts; the wider screens also used a
persona-minus-neutral direction for every persona. Candidates are ranked by
$\cos\big(\Wdown^{(L)}[:,n],\, d_{p}^{(L)}\big)$, in the first screen over layers 6--57 and afterwards over
14--48, and the dose sign is the sign of the cosine. In the first screen we also ranked by a paired $t$
statistic on each neuron's activation between the persona's answers and the other conditions' answers to the
same prompt, taking the top 20 from each list per persona.

\subsection{Screening, confirmation, controls}\label{sec:protocol}
\emph{Screen:} every candidate at dose $0.75$ with its ranked sign on the six SCREEN prompts (the prompts of
the original demonstration), 80 tokens, alongside random neurons from random layers. Finalists were chosen by
an exploratory trait-lift score on the screen.\footnote{During the run, finalists were picked by
$(m-\mu)/(\sigma+0.03)$ against the screen's controls, where $m$ is the condition's mean trait probability and
$\mu,\sigma$ the controls'. The $0.03$ floor makes that score a scaled probability difference, not a $z$
statistic, and we do not report it. Nothing in this paper's results depends on it except which neurons
reached confirmation.} \emph{Confirm:} each finalist on a prompt set it was not selected on (HOLDOUT, 12
prompts, for finalists chosen on SCREEN; HOLDOUT2, 12 fresh prompts, for later rounds), 160 tokens, at doses
$0.5$, $0.75$ and $1.0$ of the ranked sign and at $0.75$ of the reversed sign. Each confirmation set carried
random neurons drawn from the finalists' layers at the same dose as controls: 16, 13, 12 and 1 neuron in the
four rounds, and 6 for the dose-curve run (Table~\ref{tab:controls}). The unsteered model answered every set.

\subsection{Judge}\label{sec:judgemethod}
Each reply is scored once by \texttt{typesafe/jev-1.14} \citep{jev2026} on 16 yes/no questions (13 traits plus
\emph{strange}, \emph{coherent}, \emph{on\_topic}; Appendix~\ref{app:judge}), each returning a probability.
The judge is told to treat the text as data and to judge the reply's own voice. Empty replies and replies
that begin with a leaked \texttt{<think>} tag were judged like any other and are included in every mean.

\subsection{Statistics and inclusion rule}\label{sec:rule}
For a finalist at a dose, the \emph{top trait} is the trait question with the highest mean probability over
the set's prompts, and $P$ is that mean. $\Delta P$ is $P$ minus the mean probability of the same trait over
that set's control replies; its 95\% interval is a bootstrap over prompts and over control replies (4{,}000
draws). A \emph{hit} is a reply with top-trait probability at least $.5$; hit rates carry Wilson intervals.
A neuron \emph{passes} if, at the screening dose magnitude $0.75$ on a set it was not selected on, $P\ge .5$,
the $\Delta P$ interval excludes zero, and mean coherence is at least $.6$. The $.5$ floor is the hit
threshold applied to the mean; the coherence floor is the one used informally during the run. Other doses are
reported but do not decide inclusion. The rule was fixed after the run and before the analysis in this paper,
and every finalist is reported against it (Table~\ref{tab:catalogue} and Appendix~\ref{app:failed}).

\subsection{Judge validation}\label{sec:valmethod}
To check that the judge tracks a human reading, the author labeled a blind, stratified sample of 123 held-out
replies: 60 from 15 conditions at the doses an earlier draft reported (4 each), 22 random controls, 16
unsteered baseline, 12 from the layer-33 dial at $+1.0$ and $-0.5$, and 13 steered replies whose top judge
trait fell between .3 and .7. The labeler saw prompt and reply only, marked every trait that applied, and
judged coherence and on-topic yes or no. The labeler had read the lab notebook, which quotes five of these
123 replies with their conditions; the rest had not been seen. Section~\ref{sec:judge} reports agreement.

\subsection{Compute}
One Mac (Apple M4 Max, 64\,GB). 1{,}767 steering jobs covering \nScreened{} distinct candidate neurons, plus
controls and baselines, and 12{,}018 judge calls (\$0.59) between 22:37 and 03:41.
